\scanpage{15}
\textcircled{7} Réconstruction du \unsure{I|1}-groupoïde fondamental $\Pi_1(X';J)$ de $X'$, et des opérations de $\mathcal G'$ dessus, en termes de l'extension $\mathcal E'$ de $\mathcal G'$ par $\pi_1$, et de l'opération de $\mathcal E'$ sur $\widetilde X'(J)$ et $\widetilde X'(J)\longrightarrow J$ [et des applications (7) (8)
\diagram{\begin{tikzcd}[ampersand replacement=\&]
\Psi^{\unsure{\mathrm{ex}|\mathrm{ext}}}:\widetilde X'(J) \arrow[r] \& \operatorname{Splitt}(\mathcal E'/\mathcal G')/\pi_1
\end{tikzcd}}
\[
\Psi:\bigl(\diagram{\begin{tikzcd}[ampersand replacement=\&]
\widetilde X'(j) \arrow[r,dash] \& \operatorname{Splitt}(\mathcal E'_j/\mathcal G'_j)
\end{tikzcd}}\bigr)].
\]
\margin{inutile de se les donner, elles résultent des autres}
\[
\operatorname{Ob}\Pi_1=J
\]
\deleted{\unsure{Soit} $j,j'\in J$}

\uline{NB} si $\widetilde{\unsure{\xi|j}}\in\widetilde X'(J)$, $\widetilde X'[\unsure{j|\xi}]$ (rev universel de $X'$ en $\unsure{\xi|j}$) s'identifie à $\widetilde X'$ (puisque $\widetilde{\unsure{\xi|j}}$ est un pt marqué de $\widetilde X'$ au dessus de $\unsure{\xi|j}$
\diagram{\begin{tikzcd}[ampersand replacement=\&]
\widetilde X' \arrow[r,dash,"\unsure{\varphi|\Psi}_{\widetilde\xi}"] \& \widetilde X'[\unsure{j|\xi}]
\end{tikzcd}}
Si $\gamma\in\pi_1$, on aura
\[
\Psi_{\gamma\cdot\widetilde{\unsure{\xi|j}}}=\Psi_{\widetilde{\unsure{\xi|j}}}\circ\gamma_{\widetilde X'}^{\deleted{\unsure{-1}}}
\]
donc on peut écrire, indépendamment du choix du pt base $\widetilde{\unsure{\xi|j}}$
\diagram{\begin{tikzcd}[ampersand replacement=\&]
\widetilde X'[j] \& \underset{\text{torseur à dr sous }\pi_1}{\underbrace{\widetilde X'(j)}}\mathbin{\wedge}_{\pi_1}\widetilde X' \arrow[l,"\sim"']
\end{tikzcd}}
\mspage{14}
